\subsubsection*{4.9.}
回到 4.1 的较强假设，设 $F$ 是 $X$ 上的凝聚层，其在 $X^*$ 上的限制局部自由。

\subsubsection*{命题 4.10.}
\underline{设 $f$ 是 $F$ 在 $\bar X^*$ 上的逆像 $\bar F$ 的一个有限决定型截面。则 $f$ 沿 $X_0$ 适度增长，当且仅当分解 (4.7.1) 中的各 $f_{\alpha,k}$ 沿 $X_0$ 亚纯。}

“适度增长”的定义见 [1], 2.10。由 (4.7.1)、(4.7.2) 和 (4.7.3) 可清楚看出，$f$ 适度增长当且仅当各 $f_{\alpha,k}$ 适度增长。对 $X^*$ 上 $F$ 的截面，已知沿 $X_0$ 适度增长等价于沿 $X_0$ 亚纯。

对 $X=D$ 及 $F=\mathcal O$，4.10 有如下初等推论。

\subsubsection*{推论 4.11.}
\underline{设 $K$ 为 $D$ 上亚纯函数在 $0$ 处的芽构成的域，并设 $K^m$（分别为 $K^{mu}$ 和 $K^{mqu}$）为 $\widetilde D^*$ 上在 $0$ 处适度增长且为有限决定型（分别为幺幂有限决定型和拟幺幂有限决定型）的全纯函数在 $0$ 处的芽构成的 $K$-代数。则以下分别是它们在 $K$ 上的基：}
\[
\begin{array}{lll}
K^m:&z^\alpha(\log z)^k,
  &0\leq\operatorname{Re}(\alpha)<1,\ k\in\mathbb N,\\[2pt]
K^{mu}:&(\log z)^k,&k\in\mathbb N,\\[2pt]
K^{mqu}:&z^\alpha(\log z)^k,
  &\alpha\in\mathbb Q,\ 0\leq\alpha<1,\ k\in\mathbb N.
\end{array}
\]

在 4.9 的假设下，用 $\Psi_\eta^m(F)$（分别用 $\Psi_\eta^{mu}(F)$ 与 $\Psi_\eta^{mqu}(F)$）表示 $\bar i^*\bar j_*\bar F$ 的如下子层：它由 $\bar F$ 在 $X_0$ 的各点附近、沿 $X_0$ 适度增长且为有限决定型（分别为幺幂有限决定型和拟幺幂有限决定型）的截面构成。设
\[
F^m\subset j_*j^*F
\]
为 $X$ 上由 $j_*j^*F$ 中沿 $X_0$ 亚纯的截面所成的层。命题 4.10 有下述不依赖于 $z$ 和 $\log z$ 的任意选择的推论。

\subsubsection*{推论 4.12.}
\underline{自然态射}
\[
\textup{(i)}\qquad
K^m\otimes_K i^*F^m\longrightarrow\Psi_\eta^m(F),
\]
\[
\textup{(ii)}\qquad
K^{mu}\otimes_K i^*F^m\longrightarrow\Psi_\eta^{mu}(F),
\]
\[
\textup{(iii)}\qquad
K^{mqu}\otimes_K i^*F^m\longrightarrow\Psi_\eta^{mqu}(F)
\]
\underline{都是同构。}

本节的主要结果如下。

\subsubsection*{定理 4.13.}
\underline{由 (4.3.1) 诱导的自然态射}
\begin{equation}\tag{4.13.1}
\Psi_\eta^{mqu}(\Omega_X^*)
\longrightarrow R\Psi_\eta(\mathbb C),
\end{equation}
\underline{是 $D^+(X_0,\mathbb C)$ 中的同构。}

实际上，我们将证明更一般的结果 4.15 和 4.17。它们的陈述需要一些预备知识；只关心 4.13 的读者可以略过。

\subsubsection*{4.14.}
设 $V$ 是 $X^*$ 上的有限维复向量空间局部常值层（局部系统）。设
\[
\mathcal V=V\otimes_{\mathbb C}\mathcal O
\]
为相应的局部自由 $\mathcal O$-模，设 $\nabla$ 是 $\mathcal V$ 上使 $V$ 的局部截面水平的可积联络。已知 $\mathcal V$ 可延拓为 $X$ 上的凝聚层 $\widetilde{\mathcal V}$，使 $\nabla$ 相对于此延拓沿 $X_0$ \underline{正则}。还已知任意两个这样的延拓亚纯等价（[1], II, 5.7）。因此，层 $\Psi_\eta^m(\widetilde{\mathcal V})$ 及其变体 $\Psi^{mu}$ 与 $\Psi^{mqu}$ 不依赖于所选的延拓；将其记为 $\Psi_\eta^m(\mathcal V)$。同样定义
\[
\Psi_\eta^m\bigl(\Omega_{X^*}^p(V)\bigr)
=\Psi_\eta^m(\Omega_X^p\otimes\mathcal V),
\]
及其变体 $\Psi^{mu}$ 与 $\Psi^{mqu}$。

\subsubsection*{命题 4.15.}
\underline{由 (4.4.1) 诱导的自然态射}
\begin{equation}\tag{4.15.1}
\Psi_\eta^m\bigl(\Omega_{X^*}^*(V)\bigr)
\longrightarrow R\Psi_\eta(V),
\end{equation}
\underline{是 $D^+(X_0,\mathbb C)$ 中的同构。}

\subsubsection*{4.16.}
若对每个从标准圆盘 $D_0$ 到 $X$、满足 $u^{-1}(X_0)=\{0\}$ 的态射 $u$，$D_0^*$ 上局部系统 $u^*V$ 的单值变换均为幺幂（分别为拟幺幂），则称 4.14 中的局部系统 $V$ 沿 $X_0$ \underline{幺幂}（分别为\underline{拟幺幂}）。

\subsubsection*{命题 4.17.}
\underline{若 $V$ 拟幺幂，则由 (4.4.1) 诱导的自然态射}
\begin{equation}\tag{4.17.1}
\Psi_\eta^{mqu}\bigl(\Omega_{X^*}^*(V)\bigr)
\longrightarrow R\Psi_\eta(V),
\end{equation}
\underline{是 $D^+(X_0,\mathbb C)$ 中的同构。}

定理 4.13 是 4.17 中 $V=\underline{\mathbb C}$ 的特殊情形。